% BEGIN THERMAL_R03_SYNC_13B
% Formalización reunida de resultados ya expuestos en el cuerpo.
\subsection{Définitions structurelles et résultats quantitatifs}
\label{sec:ii-conclusiones-cuantitativas}

La définition par la structure conserve l’objet antérieur à la lecture
scalaire. Pour les constantes d’action, cet objet est une section orientée
avec une norme multiplicative sur le quotient local et un caractère régional.
Pour Barbero--Immirzi, c’est
l’évaluation d’une incidence orientée. Pour le mélange, c’est un transport
entre coordonnées et amplitudes. Pour Boltzmann, c’est une application entre
droites de grandeurs. Pour la gravitation, c’est une section réciproque de
l’action dont la grandeur provient de la longueur inscrite et de la lecture
radiale spécifiées dans R1--R8. Ces définitions
déterminent des opérations différentes dans la même généalogie.

Le tableau suivant fixe, pour chaque grandeur, l’objet qui la définit,
l’opération qui détermine sa lecture et la fonction physique développée
dans le corps du texte. Les preuves de reconstruction et de transport établissent
comment cette information est conservée lors de la réutilisation du résultat. En
particulier, \ref{prop:registro-36-residuos} caractérise la représentation
du registre, \ref{mem:prop-schur} conserve l’observation sous réduction
et \ref{cor:ii-witt-composed-reader} conserve les coordonnées angulaires
dans la réalisation d’incidence. Ces trois opérations ont des domaines
propres et des fonctions démonstratives distinctes au sein de la même séquence.

\begingroup
\setlength{\LTpre}{0pt}
\setlength{\LTpost}{2pt}
\begin{longtable}{@{}p{.13\linewidth}p{.415\linewidth}p{.395\linewidth}@{}}
\caption{Objet structurel, lecture produite et sens physique des
résultats de l’article.}\label{tab:ii-conclusiones-estructura}\\
\toprule
Grandeur & Définition structurelle et résultat & Lecture physique et développement interne\\
\midrule
\endfirsthead
\toprule
Grandeur & Définition structurelle et résultat & Lecture physique et développement interne\\
\midrule
\endhead
\bottomrule
\endfoot
Planck & Section de la droite d’action :
$\hbar_{\rm pre}=H_5\,10^{-34}\mathcal U_S$,
$\hbar_{\rm ret}=\eta_{\rm ret}\,10^{-34}\mathcal U_S$.
Le quotient local a seize feuillets et
$\operatorname{ord}_{10}(\varphi\alpha^{16})=-34$.
& La section relie la fréquence et l’énergie ; le tour complet détermine
$h_a=2\pi\hbar_a$. Développement dans \ref{sec:action} ; coordonnées et comparaison
dans \ref{tab:ii-planck-comparacion}.\\[5pt]
Barbero--Immirzi & Évaluation orientée des canaux et de l’incidence :
$\gamma_{\HMT}=\pi AC^*/180+12\Delta S_{90}-\Delta S_{120}$.
L’évaluation imprimée est $0.2740076726944592747\ldots$.
& Le coefficient intervient dans la résolution élémentaire d’aire et dans
la réalisation Cartan--Holst. Les définitions de
\ref{sec:barbero} et \ref{sec:area} conservent ces deux usages ;
\eqref{eq:holst} spécifie le second.\\[5pt]
Mélange CKM & Image du vecteur angulaire par les trois applications
sectorielles ; phase $\delta_{\deg}=9A_{\deg}$.
Les angles et les amplitudes sont évalués après la construction du lecteur.
& Le quartet de Jarlskog exprime l’information de phase invariante
sous les changements diagonaux de base. Les règles, leur évaluation et la
réalisation unitaire sont développées dans
\ref{sec:ii-propagacion-angular}.\\[5pt]
Boltzmann & Transduction positive
$k_B:\mathcal L_\Theta\to\mathcal L_E$ avec
$k_B\Theta_{{\rm clk},a}\log3=h_a/T_{108}$.
L’état et la mesure des continuations déterminent le contenu
informationnel sur lequel elle agit.
& $S=k_Bs$ réalise dimensionnellement l’entropie de l’état réduit.
Les bases d’énergie et de température et leurs transformations sont conservées
dans \ref{sec:ii-boltzmann} ; la caractérisation additive et l’instrument
sont démontrés dans \ref{sec:ii-aditiva-instrumento} ; la comparaison figure dans
\ref{tab:ii-kb-g-cartas}.\\[5pt]
Réalisation thermique marquée & Référence fixe $\eta$, avec les multiplicités
$(d_0,d_1,d_2)=(1,380,649)$ aux degrés $23,26,29$ et la trace
$b_*=1.380649\times10^{-23}$.
Les lectures $\mathcal S_P=b_*s$ et $\mathcal E_P=\operatorname{Tr}(\rho H)$
déterminent $d\mathcal S_P/d\mathcal E_P=b_*\beta$ et
$\Theta_P=(b_*\beta)^{-1}$.
& Avec une référence fixe, des bases positives transportées et un hamiltonien non
scalaire, \ref{thm:ii-ter27-transducer} détermine le transducteur.
Le corollaire~\ref{cor:ii-ter27-variational} prouve le minimum unique.
La composition conserve $q^N/Z_N(q)$, l’origine énergétique et les
transports horloge--horizon de \eqref{eq:ii-ter27-intertwiner}.\\[5pt]
Gravitation & Section réciproque
$G_a=c_{\rm int}^{3}L^2/\hbar_a$,
avec $L=(5759/23040)\alpha_{\HMT}^{16}L_*$.
L’horloge $t_0=\pi_{\HMT}L/(54c_{\rm int})$ conserve l’expression
circulaire $G_a=\vartheta_{108}^{\circ}c_{\rm int}^{5}t_0^2/\hbar_a$,
avec $\vartheta_{108}^{\circ}=(54/\pi_{\HMT})^2$.
& Le produit $G_a\hbar_a$ conserve la normalisation aréale.
La construction longitudinale et radiale, ses règles et l’unicité du
coefficient sont démontrées dans \ref{g26:sec:rigidez-radial-es} ;
la réalisation circulaire est conservée dans \ref{sec:gravity} et la confluence thermique dans
\ref{sec:confluencia}.\\
\end{longtable}
\endgroup

La comparaison quantitative conserve le sens propre de chaque
référence. Pour Planck, le tableau~\ref{tab:ii-planck-comparacion} compare
des coordonnées dimensionnelles à une constante définissant le SI ; les
résidus relatifs des orientations antérieure et postérieure sont,
respectivement, d’environ $-1.82\times10^{-15}$ et
$-6.93\times10^{-10}$. Pour Barbero--Immirzi, le
tableau~\ref{tab:ii-barbero-comparacion} compare deux constructions :
l’évaluation interne et la racine de l’équation de Ghosh--Mitra, dont la
valeur approchée est $0.2740668582328300$. La différence relative est
d’environ $-2.16\times10^{-4}$ ; l’objet comparé est un
coefficient théorique.

Pour CKM, le tableau~\ref{tab:ii-ckm-comparacion} réunit les sinus des
trois angles de mélange, la phase en radians et l’invariant de Jarlskog.
Les angles et la matrice des modules sont explicités dans le développement
adjacent.
Les écarts marginaux normalisés qui y sont définis satisfont
$|z|<0.029$ par rapport aux lignes PDG utilisées. Chaque résidu utilise
l’incertitude marginale indiquée. Leur lecture est individuelle ;
une évaluation conjointe nécessite en outre la matrice de covariance
de ces observables.
Pour Boltzmann et la gravitation, le tableau~\ref{tab:ii-kb-g-cartas}
expose séparément le transducteur, la composition longitudinale et radiale, les coordinatisations dimensionnelles
et la référence externe. Cette organisation permet
d’identifier exactement la quantité que produit chaque construction et
celle qu’évalue chaque comparaison.

Le tableau~\ref{tab:ii-conclusiones-contraste} réunit les résultats de
ces comparaisons avec les définitions structurelles précédentes.
\begin{table}[!htbp]
\centering\small
\setlength{\tabcolsep}{4pt}
\renewcommand{\arraystretch}{1.12}
\begin{tabularx}{\linewidth}{@{}>{\raggedright\arraybackslash}p{.16\linewidth}>{\raggedright\arraybackslash}X>{\raggedright\arraybackslash}X@{}}
\toprule
Grandeur & Résultat dans la réalisation spécifiée & Référence et comparaison\\
\midrule
Planck, avant le retour
& \(\hbar_{\rm pre}=1.0545718176461544696\ldots\)
  \(\times10^{-34}\,\mathrm{J\,s}\).
& SI : \(\hbar_{\rm SI}=h_{\rm SI}/(2\pi)\), exact.
  Sa coordonnée est \(1.0545718176461563912\ldots\)
  \(\times10^{-34}\,\mathrm{J\,s}\).
  Résidu relatif \(-1.822221\times10^{-15}\).\\[4pt]
Planck, après le retour
& \(\hbar_{\rm ret}=1.0545718169150390611\ldots\)
  \(\times10^{-34}\,\mathrm{J\,s}\).
& Même référence SI.
  Résidu relatif \(-6.932836\times10^{-10}\).
  Pour \(h_a=2\pi\hbar_a\), les résidus sont identiques.\\[4pt]
Barbero--Immirzi
& \(\gamma_{\HMT}=0.2740076726944592747\ldots\).
& Ghosh--Mitra : environ \(0.2740668582328300\).
  Différence relative \(-2.159529202\times10^{-4}\)
  entre constructions théoriques.\\[4pt]
Mélange CKM
& \(\delta_{\rm CP}=65.676173123554^\circ\),
  \(J=3.1189723326632974\times10^{-5}\).
  Les trois angles de \eqref{eq:ii-ckm-decimales} sont
  \(\theta_{12}=13.003313589509^\circ\),
  \(\theta_{23}=2.398056157332^\circ\) et
  \(\theta_{13}=0.213977382244^\circ\).
& PDG 2025 : les cinq comparaisons de
  \(s_{12},s_{23},s_{13},\delta,J\) satisfont \(|z_x|<0.029\)
  avec les incertitudes marginales du
  tableau~\ref{tab:ii-ckm-comparacion}.\\[4pt]
Boltzmann
& \(B_*=1.380649\times10^{-23}\,u_E/u_\Theta\),
  avec les bases et la règle thermométrique de
  \eqref{eq:ii-b-evaluacion-comparada}.
& SI : \(1.380649\times10^{-23}\,\mathrm{J/K}\), exact.
  Coïncidence du nombre sous le transport dimensionnel déclaré.\\[4pt]
Longueur \(L\)
& \(L=1.6162549826251263301\ldots\)
  \(\times10^{-35}\,\mathrm m\), pour \(L_*=1\,\mathrm m\).
& Composition longitudinale et radiale :
  \(G_{\rm ret}=c_{\rm int}^3L^2/\hbar_{\rm ret}\),
  avec les mêmes bases de vitesse et d’action.\\[4pt]
Gravitation
& \(G_{\rm ret}=6.6742996594140227\ldots\)
  \(\times10^{-11}\,\mathrm{m^3\,kg^{-1}\,s^{-2}}\).
& CODATA 2022 : \(6.67430(15)\times10^{-11}\) dans les mêmes unités.
  Différence d’environ \(-0.00227057\) incertitudes-types.\\
\bottomrule
\end{tabularx}
\caption{Synthèse de la comparaison quantitative. Les évaluations, les bases
et les références sont celles des tableaux~\ref{tab:ii-planck-comparacion}--\ref{tab:ii-kb-g-cartas}.
Les résidus relatifs des constantes d’action, la comparaison
théorique de Barbero--Immirzi et les résidus normalisés de CKM et de la gravitation
conservent leurs sens respectifs.}
\label{tab:ii-conclusiones-contraste}
\par\medskip
\centering\small
\setlength{\tabcolsep}{4pt}
\renewcommand{\arraystretch}{1.18}
\begin{tabular}{@{}lrrr@{}}
\toprule
Observable & Évaluation HMT & PDG 2025 & \(z_x\)\\
\midrule
\(s_{12}\) & \(0.225007404757\) & \(0.22501\pm0.00068\) & \(-0.003817\)\\
\(s_{23}\) & \(0.041841757010\) & \(0.04183^{+0.00079}_{-0.00069}\) & \(0.014882\)\\
\(s_{13}\) & \(0.003734601164\) & \(0.003732^{+0.000090}_{-0.000085}\) & \(0.028902\)\\
\(\delta\) (rad) & \(1.146265461116\) & \(1.147\pm0.026\) & \(-0.028251\)\\
\(10^5J\) & \(3.118972333\) & \(3.12^{+0.13}_{-0.12}\) & \(-0.008564\)\\
\bottomrule
\end{tabular}
\caption{Valeurs CKM et comparaison marginale réunies dans les conclusions.
Les trois angles du tableau~\ref{tab:ii-conclusiones-contraste} sont comparés
au moyen de \(s_{ij}=\sin\theta_{ij}\), avec la phase en radians et
l’invariant \(J\). On conserve les chiffres, incertitudes et résidus du
tableau~\ref{tab:ii-ckm-comparacion}, correspondant à l’ajustement à trois
générations PDG 2025. La dernière ligne exprime conjointement la valeur et
l’incertitude en unités de \(10^{-5}\). Les résidus sont des comparaisons
marginales d’observables corrélées.}
\label{tab:ii-conclusiones-ckm}
\end{table}


\Needspace{10\baselineskip}
\subsection{Conservation de l’échelle aréale et récupération de l’information de frontière}

La relation entre information et gravitation se concentre dans deux applications
complémentaires. La première conserve l’échelle aréale au moyen de
$G_a\hbar_a$ et organise ses occupations au moyen de
$\gamma_{\HMT}a_0$. La seconde conserve la résolution sectorielle au moyen de
l’observation conjointe de l’aire et de l’incidence pondérée. L’inversion
\eqref{eq:recover} retrouve les occupations à partir de l’aire et de
l’observable pondérée ; les équations
\eqref{eq:ii-respuesta-entropica} et \eqref{eq:ii-respuesta-areal}
relient leurs variations aux fluctuations de l’état.

De là provient la lecture physique centrale de l’article : une même
configuration admet une observation géométrique et une observation
informationnelle reliées par une structure récupérable. Le changement
d’orientation peut conserver l’entropie de l’état réduit et
transformer l’aire, comme l’établissent
\eqref{eq:ii-inversion-entropia} et \eqref{eq:ii-area-complementaria}.
La conservation de l’information se rapporte à l’état et à son transport ;
l’entropie d’intrication se rapporte en outre à la bipartition et à
la réduction. L’identité d’horizon de
\ref{sec:confluencia} utilise ces fonctions dans la réalisation
géométrique et thermique spécifiée.

Les identités de mémoire permettent également de préciser la portée d’une
réduction de variables. L’égalité des observations de
\eqref{mem:schur-lector} conserve les contributions du secteur éliminé
au moyen de l’opérateur effectif, de la source et du lecteur transformés.
L’entropie de l’état réduit, en revanche, s’évalue avec la bipartition
et la mesure spécifiées dans \ref{sec:ii-ecuacion-informacion}.
Les deux constructions expliquent des fonctions différentes de la mémoire :
la première conserve une lecture du transport ; la seconde quantifie
l’information accessible au sous-système. Leur articulation maintient
explicites la structure transportée et la réduction qui définit
l’observable physique.

Le progrès explicatif s’exprime donc au moyen d’un ordre précis :
structure engendrée, application de lecture, valeur obtenue et fonction
physique. L’égalité de deux valeurs s’étudie avec l’égalité
ou la différence de leurs structures ; les applications de récupération déterminent
laquelle de ces différences conserve un sens observable.
% END THERMAL_R03_SYNC_13B
